\documentclass{article}
\usepackage[T1]{fontenc}
\usepackage{lmodern}
\usepackage{graphicx}
\usepackage{amsmath,amssymb}
\usepackage[a4paper,margin=2cm]{geometry}
\usepackage[absolute,overlay]{textpos}
\setlength{\TPHorizModule}{1mm}
\setlength{\TPVertModule}{1mm}
\usepackage[indonesian]{babel}
\usepackage{hyperref}
\setlength{\emergencystretch}{2em}
% unit: Halaman 12

\begin{document}

% === Halaman 12 (llm) ===

\section*{Algoritma Pertukaran Kunci Diffie-Hellman}

\begin{enumerate}
  \item Alice membangkitkan bilangan bulat acak $a$ dan mengirim hasil perhitungan berikut kepada Bob:
  \[ A = g^a \bmod p \quad (a = \text{kunci privat Alice, } A = \text{kunci publik Alice}) \]
  \item Bob membangkitkan bilangan bulat acak $b$ dan mengirim hasil perhitungan berikut kepada Alice:
  \[ B = g^b \bmod p \quad (b = \text{kunci privat Bob, } B = \text{kunci publik Bob}) \]
  \item Alice menghitung
  \[ K = B^a \bmod p \quad (K = g^{ab} \bmod p) \]
  \item Bob menghitung
  \[ K = A^b \bmod p \quad (K = g^{ab} \bmod p) \]
\end{enumerate}

\begin{itemize}
  \item Jika perhitungan dilakukan dengan benar, maka $K = g^{ab} \bmod p$. Alice dan Bob sekarang memiliki kunci rahasia yang sama, yaitu $K$.
\end{itemize}

\end{document}
